\section*{第 \ref{ch:g7:volume-mass-density} 章——体积、质量、密度}

\begin{solution}{exo:g7:volume-mass-density:1}
$\rho = m/V$：$\rho$ 表示\omterm{def:g7:volume-mass-density:formula}{密度}，$m$ 表示样品的\omterm{def:g3:measuring-mass:mass}{质量}，$V$ 表示
其\omterm{def:g6:mass-and-volume:volume}{体积}。以克和\omterm{def:g6:mass-and-volume:volume}{立方厘米}为\omterm{def:g6:measurement-in-science:measuring}{单位}时，$\rho$ 的\omterm{def:g6:measurement-in-science:measuring}{单位}是
\unit{g/cm^3}。
\end{solution}

\begin{solution}{exo:g7:volume-mass-density:2}
$\rho = 270 \div 100 = \qty{2.7}{g/cm^3}$：铝。
\end{solution}

\begin{solution}{exo:g7:volume-mass-density:3}
$m = 9.0 \times 40 = \qty{360}{g}$，即铜的\omterm{def:g3:measuring-mass:mass}{质量}。
$V = 1930 \div 19.3 = \qty{100}{cm^3}$，即金的\omterm{def:g6:mass-and-volume:volume}{体积}——正是皇冠
案例中给出的那个数。
\end{solution}

\begin{solution}{exo:g7:volume-mass-density:4}
去皮操作去掉了量筒本身的\omterm{def:g3:measuring-mass:mass}{质量}，使\omterm{def:g3:levers-and-scales:scale}{天平}只显示
\omterm{def:g3:solids-liquids-gases:liquid}{液体}的\omterm{def:g3:measuring-mass:mass}{质量}——即“诚实零点”规则。随后将
\omterm{def:g3:levers-and-scales:scale}{天平}示数 $m$ 与\omterm{met:g6:mass-and-volume:cylinder}{弯月面}读数 $V$ 代入公式。
\end{solution}

\begin{solution}{exo:g7:volume-mass-density:5}
$\rho = 70 \div 50 = \qty{1.4}{g/cm^3}$——与蜂蜜相近。冰的
\omterm{def:g7:volume-mass-density:formula}{密度}为 $0.9$，明显更小：它在水中浮得很高。
\end{solution}

\begin{solution}{exo:g7:volume-mass-density:6}
两个\omterm{def:g6:measurement-in-science:measuring}{单位}同步换算：一升为 $1000$ \unit{cm^3}，一\omterm{def:g3:measuring-mass:units}{千克}为
$1000$ 克，因此两个千倍相互抵消——同一个
数值适用于两者。将克与升混用则跳过了这一抵消，
使所得数值错误了一千倍。
\end{solution}

\begin{solution}{exo:g7:volume-mass-density:7}
$\rho = 58 \div 5.0 = \qty{11.6}{g/cm^3}$——接近铅的
$11.3$，与金的 $19.3$ 相差甚远。结论：并非纯金；
很可能是披着金子外衣的内里含铅的冒牌货。
\end{solution}

\begin{solution}{exo:g7:volume-mass-density:8}
$100 \times 100 \times 100 = 1000000$ \unit{cm^3}——即一百万。
每\omterm{def:g6:mass-and-volume:volume}{立方厘米}按 \qty{1}{g} 计，共 \qty{1000000}{g} $=$ \qty{1000}{kg}：
即每立方米水重一吨——这正是大型水床
只宜放在底层的原因。
\end{solution}

\begin{solution}{exo:g7:volume-mass-density:9}
$V = 260 - 200 = \qty{60}{cm^3}$；
$\rho = 420 \div 60 = \qty{7.0}{g/cm^3}$——介于铝
（$2.7$）与铁（$7.9$）之间，且最接近铁而略低。“介于”
之间的数值暗示：可能是混合物或镀层冒牌货，绝非银的
$10.5$。
\end{solution}

\begin{solution}{exo:g7:volume-mass-density:10}
\omterm{def:g3:solids-liquids-gases:liquid}{液体}\omterm{def:g3:measuring-mass:mass}{质量}：$1010 - 380 = \qty{630}{g}$；\omterm{def:g6:mass-and-volume:volume}{体积}：
$\qty{75}{cL} = \qty{750}{cm^3}$；故 $\rho = 630 \div 750 =
\qty{0.84}{g/cm^3}$——是一种油。
\end{solution}

\begin{solution}{exo:g7:volume-mass-density:11}
包裹的\omterm{def:g7:volume-mass-density:formula}{密度}：$1200 \div 4000 = \qty{0.3}{g/cm^3}$——远低于
水的 $1.0$，浮标会浮得很高。铝的 $2.7$
只描述金属外壳；包裹内大部分是封闭的\omterm{def:g2:air-around-us:air}{空气}，
因此决定浮沉的是包裹本身——即总\omterm{def:g3:measuring-mass:mass}{质量}除以总\omterm{def:g6:mass-and-volume:volume}{体积}——它
才服从浮沉规律。
\end{solution}

\begin{solution}{exo:g7:volume-mass-density:12}
步骤：称出链条的\omterm{def:g3:measuring-mass:mass}{质量}（$m$）；用水位上升法测其\omterm{def:g6:mass-and-volume:volume}{体积}，
即将链条完全浸没、不留气泡（$V$）；计算 $\rho = m/V$；
若为纯铜，应约为 \qty{9.0}{g/cm^3}。需要注意的疑问：空心
链环会裹住\omterm{def:g2:air-around-us:air}{空气}而使 $V$ 偏大，造成\omterm{def:g7:volume-mass-density:formula}{密度}偏低的假象；另外，狡猾的
混合物（或镀层内芯）可能接近 $9.0$，却完全不含
纯铜。进一步的检验：重复\omterm{def:g6:measurement-in-science:measuring}{测量}\omterm{def:g6:mass-and-volume:volume}{体积}，
即先让链环充分浸水（抖掉气泡）；再增加一项独立的特征检验，
例如用\omterm{def:g1:magnets:magnet}{磁铁}检验（钢芯会立即显露），
或在车间里测定其
\omterm{def:g6:water-states:changes}{熔化}特性。
\end{solution}

\begin{solution}{pb:g7:volume-mass-density:1}
\textbf{1.} $5400 \div 2000 = \qty{2.7}{g/cm^3}$：铝。
\textbf{2.} $1800 \div 200 = \qty{9.0}{g/cm^3}$：铜。
\textbf{3.} $4520 \div 400 = \qty{11.3}{g/cm^3}$：铅——\omterm{def:g7:volume-mass-density:formula}{密度}大、
质地柔软，粗心接触有毒，需戴手套。
\textbf{4.} $8100 \div 1500 = \qty{5.4}{g/cm^3}$——介于
铝（$2.7$）与铁（$7.9$）之间：是两者的混合物，
按\omterm{def:g6:mass-and-volume:volume}{体积}大致各占一半。
\textbf{5.} $V = 540 \div 2.7 = \qty{200}{cm^3}$。
\textbf{6.} 一块立方体：$m = 9.0 \times 25 = \qty{225}{g}$。卷材的
\qty{1800}{g} 可制得 $1800 \div 225 = 8$ 块完整的立方体。
\textbf{7.} $\qty{20}{kg} = \qty{20000}{g}$；$20000 \div 225 =
88.9$：取整得 $88$ 块。
\textbf{8.} \omterm{def:g3:measuring-mass:mass}{质量} $m = 7.9 \times 3000 = \qty{23700}{g} \approx
\qty{23.7}{kg}$。
\textbf{9.} $3860 \div 200 = \qty{19.3}{g/cm^3}$——恰与
金的\omterm{def:g7:volume-mass-density:formula}{密度}吻合。
\textbf{10.} 钨——\omterm{def:g7:volume-mass-density:formula}{密度}为 $19.3$，与金的特征数值一致到小数点后，
\omterm{def:g7:volume-mass-density:formula}{密度}检验只能筛选候选物质，无法
区分恰好具有相同特征的物质。
\textbf{11.} 两台仪器都没有出错：\omterm{def:g3:measuring-mass:mass}{质量}和\omterm{def:g6:mass-and-volume:volume}{体积}均正确，
除法也正确。检验如实表明，金锭
\emph{要么}是金，要么是\omterm{def:g7:volume-mass-density:formula}{密度}与金完全相同的物质；
它只把嫌疑缩小为两种可能。
\textbf{12.} 可取的方法：敲击金条所发出的\omterm{def:g3:sound-around-us:sound}{声音}与手感
（钨硬得多，在隐蔽处做锉削或硬度测试即可迅速区分），
或公认的专业方法：\omterm{def:g6:measurement-in-science:measuring}{测量}金条导热或导电的能力
（金是最好的\omterm{def:g4:conductors-insulators:conductor}{导体}之一，钨则远远落后）——
比任何熔炉都温和，
与\omterm{def:g7:volume-mass-density:formula}{密度}结合便可作出决断。
\textbf{13.} 例如：“\omterm{def:g7:volume-mass-density:formula}{密度}能证明一种物质
\emph{并非}它所声称的那样。它只能提示它是什么——
特征可以缩小嫌疑范围，但仅凭\omterm{def:g7:volume-mass-density:formula}{密度}永远无法定罪：
认定身份需要两个相互独立的证据。”
\end{solution}
